\section{Background and Related Work}
\label{sec:related}

\subsection{Cloud assistants and the memory feature}
Commercial assistants have converged on a feature set that users now expect: persistent memory of facts about the user, searchable conversation history, projects that group chats with instructions and files, a canvas for co-editing documents and code, a code interpreter, web search with citations, file upload, image generation, voice conversation, and connectors to external tools. \hivi{}'s feature list (\sect{sec:knowledge}--\sect{sec:autonomy}) deliberately mirrors this set, because these are the capabilities that make an assistant useful day to day. The difference is not the list but where it runs and who controls it. Each cloud feature has a hidden assumption of elastic compute and vendor-held state; a large part of this paper is about removing those assumptions one feature at a time.

\subsection{Local inference}
\texttt{llama.cpp}~\cite{llamacpp} established that quantised transformer inference on CPUs and consumer GPUs is practical, and its GGUF file format and server (\texttt{llama-server}) have become a de facto standard for local deployment, including an OpenAI-compatible HTTP API, parallel decoding slots, JSON-schema-constrained decoding and prompt caching. Colibr\`i~\cite{colibri} is a dependency-free C engine aimed at large MoE models: it treats storage, RAM and VRAM as one hierarchy and streams the experts selected per token from disk, so models far larger than the GPU can run on consumer machines. \hivi{} uses both engines behind one proxy and chooses between them per model (\sect{sec:models}). The MoE architecture itself~\cite{shazeer2017,fedus2022} is what makes this possible: Qwen3's 30B-A3B class~\cite{qwen3} and gpt-oss-20b~\cite{gptoss} activate only a few billion parameters per token, so their decode cost fits an 8~GB card even though their weights do not fit it outright. Speculative and multi-token-prediction decoding~\cite{leviathan2023,deepseekv3} offer further speed-ups on the same hardware and appear in \hivi{}'s engine plan.

\subsection{Memory for assistants}
Long-term memory for LLM agents has been studied as an operating-system-like paging problem (MemGPT~\cite{memgpt}), as a reflective memory stream for simulated agents~\cite{park2023}, and as retrieval-augmented generation (RAG)~\cite{lewis2020} over a store of past interactions. Commercial memory features expose a small list of user-editable facts plus relevant recall. \hivi{} separates two stores with different cost profiles: a lightweight personal store of short facts with a profile block and per-message recall (\sect{sec:memory}), and a heavier project memory bank for agents (Hindsight, part of the existing foundation) whose LLM-based fact extraction takes on the order of a minute per save on the target hardware. The design decision that follows---extract after the reply, on an idle GPU, with de-duplication and undo---is driven directly by the single-GPU constraint.

\subsection{Agent frameworks and ``software companies''}
ReAct~\cite{react} and its descendants interleave reasoning and tool calls; MetaGPT~\cite{metagpt} and ChatDev~\cite{chatdev} organise LLM agents as a software company with CEO, CTO, programmer and tester roles and show that the structure works for small applications. Their reported weak points---verification and integration---are exactly what \hivi{}'s agent company (\sect{sec:ide}) targets with locked tests written before code, contracts fixed before splitting, code-run gates at every handoff, failure routing to the owning level, and budgets with stall and oscillation detection. SWE-bench~\cite{swebench} established hidden-test evaluation of code agents; \hivi{}'s local benchmark suite follows the same idea at the scale of one user's projects. LangGraph checkpointing, already used in the prototype's ProjectRoom, provides crash recovery.

\subsection{Tool use and connectors}
The Model Context Protocol (MCP)~\cite{mcp} standardises how assistants discover and call tools, resources and prompts offered by local or remote servers, and has been adopted by major assistants and IDEs. \hivi{} adopts it as the one way to connect to external apps (\sect{sec:capabilities}), but adds what small local models need: a tool router that exposes only the few relevant tools per message, grammar-constrained tool calls, step budgets, and an approval gate on anything that changes state.

\subsection{Security of LLM agents}
Indirect prompt injection~\cite{greshake2023}---instructions hidden in content the model reads---is the central threat for an assistant that reads web pages, files and email and can act. The ``lethal trifecta'' of private data, untrusted content and the ability to exfiltrate is well known in practitioner discussion; \hivi{} will have all three and designs its security base (\sect{sec:security}) around that fact: fetched content is tagged as data, actions that send, delete or spend require approval while such content is in the conversation, secrets never enter prompts, and untrusted code runs sandboxed. Sycophancy---models agreeing to please---is documented in preference-tuned LLMs~\cite{sharma2024}; \hivi{}'s loyalty rules (\sect{sec:soul}) explicitly reject it.

\subsection{Parameter-efficient adaptation}
LoRA~\cite{lora} and QLoRA~\cite{qlora} make fine-tuning feasible on consumer GPUs for models of a few billion parameters. \hivi{} uses them only as the last and most expensive of three layers for shaping character (\sect{sec:soul}), after retrieval over a knowledge library and preference data collected from the user, and only with a regression test set and rollback, because fine-tuning can degrade other abilities.

\subsection{Speech and images}
Whisper~\cite{whisper} and its optimised implementations provide local speech-to-text; small text-to-speech models (Kokoro, Chatterbox) and diffusion image models with 4-bit quantisation now fit beside a chat model on an 8~GB GPU, at the cost of careful GPU sharing (\sect{sec:capabilities}). Full-duplex speech models exist but are, at the time of writing, too large for the target hardware; \hivi{} builds the cascade first and watches that space.
